\section{Introduction}

Stars of roughly solar type show chromospheric activity that is cyclic in some, irregular in others and nearly constant in a few \cite{baliunas1995}. The Sun itself went through the Maunder minimum, when the cycle was strongly suppressed \cite{usoskin2017}. Whether a star that is about to lose its cycle gives a warning is the question we take up here. The idea comes from dynamical systems: near a transition where a stable state disappears, the recovery from small perturbations slows down, and that shows up as a rising lag-1 autocorrelation and rising variance in the time series \cite{scheffer2009,held2004}. Rolling-window indicators of this kind, with significance tested against surrogate series, are the standard recipe \cite{dakos2012}.

Two data sets make a test possible. The Mount Wilson HK Project followed dozens of stars for decades, and Baum et al.\ \cite{baum2022} assembled the series of 59 stars from Mount Wilson and Keck with instrument labels. Luhn et al.\ \cite{luhn2022} showed that HD~166620 went from a cyclic state to a grand minimum, which gives one star with states on both sides of the transition. The Sun has its own calibrated S-index record \cite{egeland2017}.

We ask two questions, declared before any indicator was computed on a real star. H1 (confirmatory) compares the S-index variability of HD~166620 across its three states. H2 (exploratory, one star) asks whether the rolling indicators of HD~166620 rise before the grand minimum. The temporal indicators are also computed in every other star with a long enough Mount Wilson segment; those results are descriptive, and the sample is too small to carry a group comparison on its own.
